\section{Conclusion}
We introduced \ourmethod, a pose-free, feed-forward framework that unifies reconstruction and instruction-driven editing of 3D scenes from unposed, instruction-edited images. By directly predicting edited 3D Gaussian splats in a single pass—without test-time optimization or pose estimation—our approach departs from the conventional reconstruct-edit-refit pipeline and enables fast, photorealistic, and instruction-aligned rendering. Two design choices are key to making this practical at scale: (i) a SAM2-based recoloring strategy that provides reliable, cross-view-consistent supervision despite the lack of true multi-view edited ground truth; and (ii) an asymmetric input scheme that pairs an edited reference view with unedited auxiliary views, encouraging the network to fuse disparate observations while preserving scene structure. Coupled with lightweight 3D regularization on Gaussian centers, these components yield robust multi-view coherence even when the input edits are imperfect or inconsistent.

To enable fair and reproducible comparisons, we proposed {\ourbench}, a scene-level benchmark spanning 20 diverse scenes and four edit types (Add / Remove / Modify / Global). On this benchmark, {\ourmethod} achieves stronger text-image alignment and improved 3D consistency than recent optimization-based and feed-forward baselines, while being orders of magnitude faster at inference. Ablations confirm the necessity of view-consistent recoloring, asymmetric inputs, and 3D geometric losses, and demonstrate that our method is editor-agnostic and generalizing to a range of 2D editing front-ends.

\vspace{1mm}
\noindent\textbf{Limitations and Future Work.}
Although recoloring provides stable supervision, it does not fully capture large geometric changes (\eg, substantial add/remove edits) or extreme material and illumination shifts. Looking ahead, promising directions include expanding supervision beyond recoloring with view-consistent generative augmentation, learning per-object disentangled controls for precise 3D edits, incorporating uncertainty-aware rendering to handle ambiguous or conflicting edits, extending to dynamic scenes and longer view sequences, and further enriching {\ourbench} with more diverse and complex scenes.


% Although recoloring offers stable supervision, it underrepresents large geometric changes (\eg, substantial add / remove edits) and extreme material / illumination shifts. Future directions include: 
% scaling supervision beyond recoloring via view-consistent generative augmentation; learning per-object, disentangled controls for precise 3D edits; 
% incorporating uncertainty-aware rendering to handle ambiguous or conflicting edits; 
% extending to dynamic scenes and longer view sequences; 
% and broadening {\ourbench} with richer scenes, stronger metrics, and human studies.